

\chapter{Thresholds}\label{ssec:threshold_theorem}

The cornerstone of QEC codes and their usefulness is based on the threshold theorem. 
It states that, even under the assumption that all components are noisy,
we can still arbitrary reduce the logical error rate $p_{log}$ of the encoded qubit, 
given that the error rates of the noisy components $p_{phy}$ is below a certain threshold $p_{th}$ if we use a fault tolerant QEC code \cite{gottesmann_book,topological_quantum_memory}.
In short, if the error rate is below this threshold the QEC code can suppress errors faster than they occur.

We can reduce the logical error rate for a fixed physical noise rate by increasing the distance $d$ of the code if the fixed physical noise rate is below the threshold.
The overhead in the number of qubits has been shown to be polynomial for the surface code \cite{topological_quantum_memory}.
Note that, while below the threshold each additional qubit reduces the logical error probability, above the threshold the logical probability rises with increased code distance.

Based on this theorem we do not need to realize perfect qubits in an experiment to enable arbitrarily long quantum computations.
Instead we can use QEC codes to suppress the logical error rates to levels, where long quantum computations become feasible.  
This makes quantum computational possible, as we do not need to create perfect qubits, just decent enough qubits to enable quantum computing, if we utilize QEC codes.

The threshold is a key characteristic of a QEC setup.
It describes the physical noise rate at which scaling up the QEC setup becomes advantageous.
Codes with higher threshold can be used with less optimal qubits and might therefore be advantageous.
% We will determine the threshold for the surface code with different choices of error model, syndrome extraction circuit and decoder in this thesis.
% \textcolor{red}{Put ref to result here}.

There are different possible approaches to determine the threshold of a QEC setup.
It can be determined through a mapping to statistical physics models for some QEC setup, including the surface code under code capacity noise\footnote{Note that under code capacity noise the syndrome extraction circuit is neglectable.}.
For those models, we can identify a phase transition between an ordered and unordered phase at the threshold \cite{fundamental_thresholds,topological_quantum_memory}.

In this thesis we will use a numerical approach to determine the threshold for different configurations as described in the following section.

\section{Numerical Determination of Thresholds}

To numerical determine the threshold we first simulate the quantum memory experiment.
In this thesis we keep the encoding (see section \ref{sec:kitaev_surface_code}) constant 
and vary the choice of error model (section \ref{ssec:error_models}), syndrome extraction circuit (section \ref{ssec:syndrome_extraction}), decoder (section \ref{sec:decoding_problem_general}) and the logical observable.
For each of those QEC configuration we calculate the logical error rate $p_{log}$ for different physical error rates $p_{phy}$ and code distances $d$.

We calculate the logical error rate $p_{log,d}(p_{phy})$ by simulating the quantum circuit.
The simulation start with the data qubit in a known state dependent on the observable we want to measure.
If we want to measure the logical $\bar{Z}$ ($\bar{X}$) observable we encode the data qubit at the start of our experiment in the logical encoded $\ket{0}_L$  ($\ket{+}_L$) state.
We then simulate the circuit with the given error model and syndrome extraction circuit to generate measurement results for the logical observable and syndromes, using the quantum simulator \texttt{stim} \cite{stim}.
In the last step we use the choosen decoder, to decode the syndrome returning a prediction for the occurred error.
Utilizing pauli frame tracking (see section \ref{sec:pauli_frame_tracking}) we can compare the predicted measurement with the result of the simulation. 
We count the number of times the prediction disagrees with the simulation result $n_{errors}$ to calculate the logical error rate for a given number of shots $n_{shots}$ 
\begin{equation}
    p_{log} = n_{errors}/n_{shots}.
\end{equation}
The uncertainty on the logical error rate $\delta_{p_{log}}$ is then
\begin{equation}
    \delta_{p_{log}} = \sqrt{\frac{p_{log}(1-p_{log})}{n_{shots}}}.
\end{equation}

Note that in this approach we only detect $X$-type or $Z$-type logical errors, dependent on the choice of logical observable, $\bar{Z}$ or $\bar{X}$ respectively.
We therefore also only decode for one of those error types at a time and can only predict the logical error rate for one of those errors.
This is due to our choice of initial state.
We can not detect if an logical $X$-type ($Z$-type) error occurred if we use a $\ket{+}_L$ ($\ket{0}_L$) state, given that those error operators are eigenstate of those logical operators.
\textcolor{green}{Is this following line useful? it is interesting to know that it is possible, but once again i did not do that.}
One could realize a simulation detecting both type of errors at the same time using a reference system to create an initial pure state, as used in \cite{quantum_data_processing,fundamental_thresholds} and referred to as coherent information setup.


An example of how the logical error curves $p_{log,d}(p_{phy})$ look like can be seen in figure \ref{}.
\textcolor{red}{Include graphic here? (or should i refer to result part)}

For the asymptotic behavior of the logical error rate, for small physical noise rates, we can state the following:
The fault tolerant distance $d$ code is guaranteed to correct $t=(d-1)/2$ independent faults, each with a noise rate proportional to the physical noise rate $p_{phy}$.
Therefore we require $t+1$ fault to cause a logical error, resulting in a logical error rate which scales in leading order as 
\begin{equation}
    \lim_{p_{phy} \rightarrow 0} p_{log} \propto  p_{phy}^{(t+1)}= p_{phy}^{((d+1)/2)}.
\end{equation}

% Necessary lines?
Above the threshold $p_{phy}>p_{th}$ each increase in the code distance $d$ leads to an increased logical noise rate $p_{log}$.
Below the threshold $p_{phy}<p_{th}$ the threshold theorem predicts that with increased code distances the logical error rate is increasingly suppressed.
At the threshold the logical error rate curves for different distances cross.

The threshold correspond to the critical value in the physical noise rate of the first-order phase transition, as shown in \cite{fundamental_thresholds}.
Therefore we can use methods from finite-size scaling analysis to identify the value of the threshold. 

\subsection{Data Collapse Method}

We use the \emph{Data Collapse Method} (DCM) to determine the threshold, utilizing the logical error rate curves $p_{log,d}(p_{phy})$ for different code distances $d$ that we simulated \cite{pyfssa}.

We can identify the system size of our problem as the code distance $d$.
The system, which is parameterized by the physical error rate $p_{phy}$, undergoes a phase transition at the critical noise rate corresponding to the threshold $p_{th}$.
For our case the order parameter is the logical error rate $p_{log}$. 
Utilizing the fact that the phase transition is a phase transition of first order, we can adapt the DCM to 
\begin{equation}
    p_{log,d}(p_{phy}) = f\left( d^{1/\nu} \cdot(p_{phy}-p_{th})\right),
    \label{eq:data_collapse_method}
\end{equation}
where $f(x)$ is the dimensionless scaling function for finite sizes and $\nu$ is the critical exponent, associated with the phase transition \cite{pyfssa}. 

Following equation \ref{eq:data_collapse_method}, plotting $y=p_{log,d}(p_{phy})$ against $x=d^{1/\nu} \cdot(p_{phy}-p_{th})$ should collapse the simulated data onto one curve, corresponding to the scaling function $f(x)$, 
under the assumption that we use the correct threshold $p_{th}$ and critical exponent $\nu$.

We can turn this relation around and determine the values of threshold $p_{th}$ and the critical exponent $\nu$, using that all data points should collapse under the correct choice of those values.
To realize this approach, we first estimate values for the critical exponent $\tilde{\nu}$ and threshold $\tilde{p}_th$.
We then use those estimates to calculate the values we want to plot against one another, which we will refer to by $x$ and $y$, using 
\begin{align}
    x&=d^{1/\tilde{\nu}} \cdot(p_{phy}-\tilde{p}_{th}) 
    \label{eq:data_collapse_x}
    \\
    y&=p_{log,d}(p_{phy})
\end{align}
We sort those data points by their $x$-values and index them as $x_i$ and $y_i$. 
We calculate a weighted difference $\Delta y_i$ between the data points and its expect value $\bar{y}_i$, which is given by linear interpolation of its neighboring data points (eq. \ref{eq:dcm_bar_y}).
We use the uncertainty on the logical error rate $\delta_{p_{log}}=\delta_y$ to compute a weighted difference between each data point $y_i$ and its expected value $\bar{y}_i$ (eq. \ref{eq:dcm_error_propagation}\& eq. \ref{eq:dcm_weighted_difference}). 
\begin{align}
\bar{y} &= \frac{(x_{i+1} - x_i)\,y_{i-1} - (x_{i-1} - x_i)\,y_{i+1}}{x_{i+1} - x_{i-1}} \label{eq:dcm_bar_y}\\
\delta_{\Delta y_i}^2 &= \delta_{y_i}^2 + \left(\frac{\delta_{y_{i-1}} (x_{i+1} - x_i)}{x_{i+1} - x_{i-1}}\right)^2 + \left(\frac{\delta_{y_{i+1}}(x_{i-1} - x_i)}{x_{i+1} - x_{i-1}}\right)^2 \label{eq:dcm_error_propagation} \\
\Delta y_i&= \frac{(y_i - \bar{y}_i)^2}{\delta_{\Delta y_{i}}^2} \label{eq:dcm_weighted_difference}
\end{align}
By minimizing the weighted difference $\Delta y_i$ all data points should fall onto a curve.

We can task the optimizer to find the minimal value for the mean of all weighted differences $\Delta y_i$, by optimizing the threshold $p_{th}$ and critical exponent $\nu$.
The optimizer we utilize is \texttt{scipy.optimize.minimize} from the \texttt{scipy} package \cite{Scipy}.
The decoder repeats the steps described above, recalculating the equations with new estimates of the threshold and critical exponent, until it minimizes the values within the given tolerance\footnote{As technical detail we will use the Nelder-Mead method with a tolerance of $10^{-5}$ \cite{Scipy}.}.
It then returns the values for the threshold $p_{th}$ and the critical exponent $\nu$, which have the minimal mean weighted difference between data points and their interpolation. 
We can then use those values to see how well such an optimization process, collapsed the data onto a curve, as shown in figure \ref{}.
\textcolor{red}{Put figure here!}

To calculate an estimated error for the threshold we use bootstrapping:
We randomly resample each determined logical error rate selecting randomly a logical error rate from its normal distribution, defined by the logical error rate $p_{log}$ and its associated uncertainty$\delta_{p_{log}}$.
Following this we redo the whole optimization process with this resampled logical error rate as input.
This process is repeated multiple times, obtaining a distribution of thresholds.
Lastly we identify the width of the central $95\%$ of this threshold distribution, as the two times the estimated error of the determined threshold.
\textcolor{green}{Should I go into more detail/resemble the line above as an equation? or should i show such an image?}
Using this method we can quantify how the uncertainty of the logical error rates propagated onto the threshold we determine for of the QEC code setup. 

The code used to implement the DCM and for bootstrap the uncertainty was adapted from code kindly provided by my one of my supervisors, Luis Colmenarez.

With the above described approach, we can determine the threshold $p_{th}$ and an associated error, given the logical error rate $p_{log,d}(p_{phy})$ for different code distances $d$.
\textcolor{red}{Maybe go into more detail, either derive the dcm method from the finite scaling ansatz or show how the data points are weighted}




